\documentclass[11pt,twocolumn]{article}
\usepackage{shared/hanzocover}
\usepackage[utf8]{inputenc}
\usepackage{amsmath,amssymb}
\usepackage{graphicx}
\usepackage{booktabs}
\usepackage{hyperref}
\usepackage{xcolor}
\usepackage{listings}
\input{shared/lstlang}

\title{The Hanzo Cloud Data Plane: Benchmarking \texttt{zip}, ZAP, and the
Road to a Million Requests per Second}

\author{
    Zach Kelling\thanks{Corresponding author: research@hanzo.ai} \\
    \textit{Hanzo Industries Inc (Techstars '17)} \\
    Los Angeles, California \\
    \texttt{https://hanzo.ai}
}

\date{2026}

\begin{document}

\hanzocoverpage

\begin{abstract}
We present a reproducible benchmark of the Hanzo cloud data plane --- the
\texttt{zip} web framework and its native ZAP binary transport --- and a
systems-tuning study on a two-node 2.5\,GbE testbed (an NVIDIA GB10 ``spark''
server, 20 cores, aarch64; a 32-core x86\_64 ``evo'' load generator). We
establish that \texttt{zip} carries a negligible tax over the \texttt{fasthttp}
engine it wraps ($\approx$24\,ns and one 48-byte allocation per request;
$0.95\times$ raw throughput on a no-op handler). We then show that the naive
distributed result --- 194{,}000 req/s over the wire --- is not a CPU or
bandwidth limit but a \emph{single hardware RX-queue} on the 2.5\,GbE NIC, whose
driver refuses hardware multi-queue. Enabling Receive Packet Steering (RPS) plus
a deeper backlog and larger socket buffers lifts the same link to
\textbf{518{,}000 req/s (a $2.67\times$ gain)}, at which point the tuned network
result exceeds loopback because the server's cores no longer contend with the
load generator. Finally we characterize the ZAP transport. It initially ran
\emph{slower} than HTTP (91k vs.\ 481k req/s loopback), which we localize by
construction to a non-zero-allocation codec --- three allocation sites, the
largest an un-pooled request context, driving 2203 garbage collections in six
seconds (relaxing \texttt{GOGC} $100\to800$ alone more than doubles throughput).
Pooling the context and making the marshalers zero-allocation --- a pure codec
rewrite with a byte-identical wire --- lifts ZAP $6.3\times$ to \textbf{570k
req/s, overtaking HTTP by $1.18\times$}, with GC falling $2203\to62$. We show
this is a CPU/allocation win, not a wire win (the ZAP frame is measurably
\emph{larger}: 124 vs.\ 75 bytes), and that over the saturated 2.5\,GbE link both
transports tie at $\approx$530k. We release the framework as
\texttt{hanzoai/benchmarks} and the tuning as a single idempotent script.
\end{abstract}

\section{Introduction}
The Hanzo control plane ships as a single Go binary that mounts every subsystem
(AI gateway, commerce, IAM, KMS) and serves \texttt{api.hanzo.ai}~\cite{hip0106}.
Its HTTP surface is \texttt{zip}~\cite{zip}, a framework built on a fork of
\texttt{fiber}~/ \texttt{fasthttp}, and its inter-subsystem calls travel over
\emph{ZAP}, a typed binary transport that reuses \texttt{fasthttp} request and
response objects over a length-framed wire~\cite{zaphttp}. Two questions follow
naturally: \emph{(i)} how much does the framework cost over the bare engine, and
\emph{(ii)} where is the real ceiling of the data plane on modern hardware ---
is it the framework, the CPU, the network stack, or the link?

Benchmarking answers only when it is fair. We hold the handler, the payload, the
concurrency, and the tuning identical across every measurement; we report
allocations, not just throughput; and we separate the \emph{load generator's}
limits from the \emph{server's}. All artifacts, scripts, and raw numbers are
public in \texttt{hanzoai/benchmarks}.

\section{Methodology}
\paragraph{Testbed.} \emph{spark} is an NVIDIA GB10 (Grace--Blackwell) developer
system: 20 aarch64 cores, running the server. \emph{evo} is a 32-core x86\_64
box, running the load generator. They are joined by a direct 2.5\,GbE wired link
(\texttt{enP7s7}$\leftrightarrow$\texttt{eno1}, both \texttt{Speed: 2500Mb/s},
subnet 192.168.77.0/24) --- confirmed with \texttt{ethtool}, distinct from the
WiFi path used for control SSH. A second host, \emph{dbc}, is wired on a 1\,GbE
USB-ethernet port for future dual-loader aggregation.

\paragraph{Framework.} \texttt{hanzoai/benchmarks} is two-tier. The
\emph{property/micro} tier (\texttt{cloud/shard}, \texttt{framework/zip}) is pure
Go, deterministic, and cluster-free, so it runs in CI and isolates a primitive
(routing ns/op) or proves an invariant. The \emph{integration} tier
(\texttt{k3s/}, \texttt{distributed.sh}) stands the real binary up and drives it
over the network. Load is generated with \texttt{bombardier} (a
\texttt{fasthttp}-based client that actually saturates the server; \texttt{hey},
a \texttt{net/http} client, caps near 100k and is unsuitable) and, for ZAP, with
a purpose-built client (\S\ref{sec:zap}) since no HTTP tool speaks the protocol.

\paragraph{Fairness.} Every server exposes the identical handler --- \texttt{GET
/health} returning the two-byte body \texttt{"ok"}. Connections are warm
(keep-alive) so we measure steady-state, not handshakes. We report the median of
sustained throughput and the p50/p99 latency; where garbage collection is
implicated we report allocations per operation from \texttt{-benchmem}.

\section{The Framework Tax: \texttt{zip} vs.\ \texttt{fasthttp}}
The in-process cost of \texttt{zip} over raw \texttt{fiber} is a constant
$\approx$24\,ns and exactly one 48-byte allocation (the \texttt{*zip.Ctx}),
independent of route shape~\cite{zipbench}; on a handler doing $\sim$3500\,ns of
real work it is under 1\%. End to end, driven by \texttt{bombardier} on loopback
(Table~\ref{tab:framework}), \texttt{zip} sustains $0.95\times$ of raw
\texttt{fasthttp} --- the wrapper is noise.

\begin{table}[t]
\centering
\caption{End-to-end HTTP, loopback, spark, \texttt{-c 125}. Same handler.}
\label{tab:framework}
\begin{tabular}{lrr}
\toprule
framework & req/s & vs.\ fasthttp \\
\midrule
raw \texttt{fasthttp} & 365--400k & $1.00\times$ \\
\texttt{zip} (HTTP)   & 338--348k & $0.95\times$ \\
\bottomrule
\end{tabular}
\end{table}

\section{The Network Ceiling Is Not the Box}
\label{sec:ceiling}
Moving the load off-box \emph{lowered} throughput --- and adding connections made
it worse (Table~\ref{tab:ceiling}). Throughput fell from 194k at \texttt{-c 1000}
to 154k at \texttt{-c 8000} while p99 latency rose $5\to52$\,ms: classic
congestion collapse. Yet the link ran at $\sim$70\,MB/s, only 22\% of the
2.5\,GbE bandwidth, and spark's cores sat mostly idle. By Little's law the
194k/5.1\,ms operating point is latency-bound, not bandwidth-bound.

The root cause is structural: \texttt{ethtool -l} reports the NIC configured with
a \emph{single hardware RX queue}, and the driver rejects reconfiguration
(\texttt{ethtool -L}: ``Operation not supported''). All inbound packets are
therefore serviced by one core's softirq, and that single core is the ceiling.

\begin{table}[t]
\centering
\caption{Untuned evo$\to$spark over 2.5\,GbE. Congestion collapse.}
\label{tab:ceiling}
\begin{tabular}{lrr}
\toprule
loader $\to$ target & req/s & p99 latency \\
\midrule
loopback (\texttt{-c 500})   & 365k & --- \\
evo$\to$spark \texttt{-c 1000} & 194k & 5.1\,ms \\
evo$\to$spark \texttt{-c 4000} & 167k & 24\,ms \\
evo$\to$spark \texttt{-c 8000} & 154k & 52\,ms \\
\bottomrule
\end{tabular}
\end{table}

\section{System Tuning: $2.67\times$ on the Same Wire}
The software remedy for a single-queue NIC is Receive Packet Steering, which
hashes flows across cores in the kernel. We apply RPS (and RFS/XPS) across all 20
cores, raise \texttt{net.core.netdev\_max\_backlog} from 1000 to 300000, set
256\,MiB socket buffers, widen the ephemeral port range, and pin the performance
governor --- all in one idempotent \texttt{tune.sh}, run on server and loader
alike. The result (Table~\ref{tab:tuned}) is $2.67\times$: the same 2.5\,GbE link
now sustains \textbf{518k req/s}, peaking at 696k.

\begin{table}[t]
\centering
\caption{evo$\to$spark, both ends tuned. Same link, same handler.}
\label{tab:tuned}
\begin{tabular}{lrrr}
\toprule
concurrency & untuned & tuned & gain \\
\midrule
\texttt{-c 1000} & 194k & 518k & $2.67\times$ \\
\texttt{-c 2000} & collapsed & 518k & --- \\
\texttt{-c 4000} & 154k & 474k (peak 696k) & $3.1\times$ \\
\bottomrule
\end{tabular}
\end{table}

Two observations. First, the tuned \emph{network} figure (518k) now exceeds the
loopback figure (365--400k): over the wire spark's 20 cores serve full-time,
whereas on loopback they are shared with the load generator. At 518k spark idles
only $\sim$12\%, so this is close to the true single-loader serving ceiling; the
remaining headroom is the province of a second, parallel NIC path.

\section{The ZAP Transport: A Zero-Alloc Codec Overtakes HTTP}
\label{sec:zap}
\texttt{zip} is ZAP-native --- a bare listen address binds the binary transport;
\texttt{http://} binds \texttt{fasthttp}. Because HTTP tools cannot speak ZAP we
built \texttt{zapbench}, a concurrent client over \texttt{zaphttp.Transport}. The
first result inverted the expectation: ZAP delivered only 91k req/s on loopback
against HTTP's 481k --- ZAP \emph{lost} by $5\times$.

The cause was not the wire but the codec. The tail (p50 303\,\textmu s but p99
13\,ms) is a garbage-collection signature, and relaxing \texttt{GOGC} $100\to800$
more than doubled ZAP to 203k --- proof the path was allocation-bound. Three
allocation sites carried it. \texttt{MarshalRequest}/\texttt{MarshalResponse}
minted a fresh frame slice per call; header decode performed \texttt{string()}
copies into a \texttt{map[string][]string}; and structurally, \texttt{zaphttp}'s
server allocated a fresh \texttt{fasthttp.RequestCtx} \emph{per request}
(\texttt{server.go}), where \texttt{fasthttp}'s own server pools it per
connection. Over a six-second run the ZAP server triggered \textbf{2203}
garbage collections; the HTTP server, 57.

The fix (\texttt{zap-proto/http}, branch \texttt{perf/zero-alloc-codec}) is
mechanical, not architectural: pool the \texttt{RequestCtx} once per connection,
hand-roll a frame encoder that appends into a pooled buffer, and drain the
4-byte length prefix through \texttt{bufio.Peek}/\texttt{Discard} so no header
array escapes to the heap. The marshalers drop to \textbf{zero} allocations
(\texttt{MarshalResponse} $19\to0$ allocs, $1215\to61$\,ns); a full request/
response round trip allocates nothing. Golden tests confirm the wire is
\emph{byte-identical} before and after --- this is a pure codec rewrite, no
protocol change. The measured effect (Table~\ref{tab:zap}): ZAP rises
$91\text{k}\to570\text{k}$ req/s, a \textbf{6.3$\times$} lift that
\textbf{overtakes HTTP at 481k --- ZAP wins $1.18\times$} on loopback, with p99
collapsing $13\,\text{ms}\to1.19\,\text{ms}$ and GC falling $2203\to62$.

Two corrections to intuition survive the fix, and both matter for where ZAP is
actually worth deploying. \emph{First, the win is CPU, not the wire.} Measured by
\texttt{TestWireBytes}, the ZAP frame is \emph{larger} than HTTP (124 vs.\ 75
bytes for a \texttt{/health} request); ZAP overtakes HTTP by spending fewer
cycles per byte, not fewer bytes. \emph{Second, \texttt{zaphttp} does not skip
JSON.} It carries the HTTP body as opaque bytes on both transports, so a handler
pays the same (de)serialization cost either way; on a chat-shaped JSON payload
ZAP is in fact $0.95\times$ HTTP, a slight loss from copying the body into the
frame tail. JSON-elision is a property of \emph{native} ZAP typed RPC --- a
separate axis \texttt{zaphttp} does not exercise. The loopback win is therefore
real and purely a CPU/allocation win; over the 2.5\,GbE link both transports tie
at $\approx$530k, saturating the wire (\S\ref{sec:wire}) where the codec is no
longer the bottleneck.

\begin{table}[t]
\centering
\caption{ZAP vs.\ HTTP, loopback, \texttt{-c 125}, same handler. The zero-alloc
codec flips ZAP from a $5\times$ loss to a $1.18\times$ win.}
\label{tab:zap}
\begin{tabular}{lrl}
\toprule
transport & req/s & note \\
\midrule
\texttt{zip} HTTP (fasthttp)          & 481k & zero-alloc; 57 GC \\
\texttt{zip} ZAP (\texttt{GOGC=100})  & 91k  & alloc-bound; 2203 GC \\
\texttt{zip} ZAP (\texttt{GOGC=800})  & 203k & $2.2\times$; confirms cause \\
\textbf{\texttt{zip} ZAP zero-alloc}  & \textbf{570k} & \textbf{0 allocs/req; 62 GC; $1.18\times$ HTTP} \\
\bottomrule
\end{tabular}
\end{table}

\section{Aside: Horizontal-Scale Routing}
The same data plane shards per-tenant SQLite state across pods by
highest-random-weight (rendezvous) election over a static membership set. A
million synthetic tenants distribute across pods with 0.05\% standard deviation
at $N{=}3$ (no hot shard), every tenant resolves to exactly one deterministic
owner (the no-two-writers invariant, verified over $10^5$ tenants $\times$ 6
re-elections), a $4\to5$ scale-up remaps only 20.1\% of tenants (the rendezvous
ideal), and the routing itself costs $\approx$545\,ns per request. Horizontal
writer scale is thus gated on I/O, not routing.

\section{Conclusions and Future Work}
The Hanzo data plane is not framework-bound (\texttt{zip} $\approx$
\texttt{fasthttp}), and on a single 2.5\,GbE link it serves 518k req/s once the
kernel's receive path is tuned --- a $2.67\times$ gain from software steering
alone, with the box only $\sim$12\% idle. The ZAP codec is now
zero-allocation and overtakes HTTP by $1.18\times$ on loopback (\S\ref{sec:zap});
the remaining levers to seven figures are orthogonal and physical: a
\emph{dual-loader} across the two physical NICs (parallel RX paths), faster
interconnect (10\,GbE/RDMA), and a \texttt{writev} split so ZAP's frame tail
avoids one body copy on large payloads. We also intend to benchmark the
post-quantum transports natively available in Go 1.26
(\texttt{X25519MLKEM768}): PQ-TLS 1.3, ZAP-over-PQ-TLS, and PQ-QUIC, separating
handshake cost (ML-KEM) from steady-state AEAD throughput.

\begin{thebibliography}{9}
\bibitem{hip0106} Hanzo Industries. \textit{HIP-0106: The Unified Cloud Binary}. 2026.
\bibitem{zip} \texttt{zap-proto/zip} --- the ZAP-native web framework. \url{https://github.com/zap-proto/zip}.
\bibitem{zaphttp} \texttt{zap-proto/http} --- ZAP-HTTP binary transport. \url{https://github.com/zap-proto/http}.
\bibitem{zipbench} \texttt{zap-proto/zip} \texttt{BENCHMARKS.md} --- framework wrapper tax.
\end{thebibliography}

\end{document}
